%! AP Statistics — Unit 4, Lesson 4.5 — Warm-Up KEY
\documentclass[10pt]{article}
\usepackage{apstats-article}
\usepackage{apstats-key}

\begin{document}

\pageheader{Unit 4, Lesson 4.5}{Warm-Up --- Answer Key}
\begin{tabularx}{\linewidth}{X X X}
Name:\;\ans{Key} & Date:\; & Period:\;
\end{tabularx}

\vspace{0.3in}

{\large\bfseries\color{navy} Part 1 \quad Spiral Review}

\vspace{0.12in}

\begin{enumerate}

    \item
    \vspace{6pt}
    \renewcommand{\arraystretch}{1.4}
    \begin{tabularx}{\linewidth}{@{} l X X r @{}}
     & \textbf{Trains Daily} & \textbf{Trains Weekly} & \textbf{Total} \\
    \hline
    \textbf{Runs}     & 24 & 16 & 40 \\
    \textbf{Swims}    & 18 & 22 & 40 \\
    \textbf{Cycles}   & 20 & 20 & 40 \\
    \hline
    \textbf{Total}    & 62 & 58 & 120 \\
    \end{tabularx}
    \vspace{6pt}
    \begin{enumerate}[label=(\alph*), itemsep=6pt]
        \item \ans{$P(\text{Runs}) = 40/120 \approx 0.333$}
        \item \ans{$P(\text{Trains Daily}) = 62/120 \approx 0.517$}
        \item \ans{$P(\text{Runs} \cap \text{Trains Daily}) = 24/120 = 0.2$}
        \item \ans{Yes, mutually exclusive. An athlete is recorded in exactly one sport row; no athlete can appear in both Runs and Swims simultaneously.}
    \end{enumerate}

    \vspace{0.18in}

    \item \ans{$0$}

    \vspace{0.14in}

    \item
    \begin{enumerate}[label=(\alph*), itemsep=6pt]
        \item \ans{$P(\text{King}) = 4/52 = 1/13$}
        \item \ans{$P(\text{King} \cap \text{Heart}) = 1/52$}
        \item \ans{13 Hearts remain to consider; 1 of them is a King.}
    \end{enumerate}

\end{enumerate}

\vspace{0.3in}

{\large\bfseries\color{navy} Part 2 \quad Looking Ahead}

\vspace{0.12in}

\noindent Today we ask: \textit{``If we already know one event occurred, how does that
change the probability of another?''}

\vspace{0.14in}

\begin{tcolorbox}[colback=greenbg, colframe=greenacc, boxrule=0.8pt, arc=2mm,
    left=4mm, right=4mm, top=3mm, bottom=3mm]
    \small
    \begin{itemize}[leftmargin=1.3em, itemsep=8pt]
      \item Conditional probability restricts the sample space to the ``given'' event.
      \item The formula is $P(A|B) = P(A \cap B) / P(B)$.
      \item From a table: numerator = joint cell; denominator = conditioning row/column total.
    \end{itemize}
\end{tcolorbox}

\vspace{0.2in}

\noindent\textcolor{slate}{\small Student responses will vary.}

\end{document}
